\section{Conclusions and Outlook}
\label{sec:conclusions}

This analysis establishes a coherent GPR-based framework for the HPS narrow prompt
$A^\prime \rightarrow e^+e^-$ resonance search. The essential ingredients are now in
place: a global GP description of the invariant-mass spectrum tied to the detector
resolution, a local profiled-likelihood extraction with an explicitly full-range signal
template, a conversion from signal yield to $\eps^2$ grounded in the radiative trident
interpretation of fixed-target heavy-photon production, and a true shared-likelihood
combination across statistically independent data sets. The accompanying toy program
already demonstrates that these ingredients can be tested in the language that matters
for a HEP search analysis, namely estimator bias, pull calibration, interval coverage,
significance calibration, and combined-search behavior.

This is a meaningful milestone for the use of Gaussian-process methods in HEP analyses.
In many applications, GPR is introduced only as a flexible interpolation tool. Here it
has been developed into a complete search framework: the GP background model is not an
isolated fit, but part of a statistically explicit inference chain that supports local
significance, expected and observed limits, and combined coupling limits across
heterogeneous data sets. That is precisely the level at which GPR must operate if it is
to become a mature analysis technique rather than an interesting auxiliary method.

The remaining work is well defined. The 2021 stages need their final production plots
and normalization inputs, the global-significance treatment needs the dedicated
look-elsewhere rerun, and the signal-absorption question should be revisited after the
next round of full procedural toys and likelihood refinements. With those steps
completed, the framework documented here is positioned to provide the first fully
combined HPS GPR-based prompt-resonance result.
